\section{Token Frequency and Variation}\label{sec:frequency}

Three frequency views are produced for every analysis, and they answer different
questions (\req{FR-17}).

\subsection{Raw versus canonical forms}

The left table is what people actually said or typed --- accents, casing and spelling
variants included. The right table collapses casing and Unicode normalisation differences.
Reading them side by side measures orthographic variation directly: a canonical form
reached by several distinct surface spellings is a documented variant cluster.

\begin{center}
\begin{minipage}[t]{0.47\linewidth}
\centering\begin{tabular}{@{}lr@{}}
\toprule
\textbf{Raw form} & \textbf{Count} \\
\midrule
\texttt{,} & 20 \\
\texttt{le} & 15 \\
\texttt{au} & 10 \\
\texttt{Le} & 9 \\
\texttt{est} & 6 \\
\texttt{je} & 6 \\
\texttt{la} & 6 \\
\texttt{pas} & 6 \\
\texttt{kwatt} & 5 \\
\texttt{on} & 4 \\
\texttt{va} & 4 \\
\texttt{Combi} & 3 \\
\texttt{a} & 3 \\
\texttt{go} & 3 \\
\texttt{les} & 3 \\
\texttt{ne} & 3 \\
\texttt{piol} & 3 \\
\texttt{à} & 3 \\
\texttt{Ashia} & 2 \\
\texttt{Je} & 2 \\
\bottomrule
\end{tabular}

\end{minipage}\hfill
\begin{minipage}[t]{0.47\linewidth}
\centering\begin{tabular}{@{}lr@{}}
\toprule
\textbf{Canonical form} & \textbf{Count} \\
\midrule
\texttt{le} & 24 \\
\texttt{,} & 20 \\
\texttt{au} & 10 \\
\texttt{je} & 8 \\
\texttt{est} & 6 \\
\texttt{la} & 6 \\
\texttt{on} & 6 \\
\texttt{pas} & 6 \\
\texttt{kwatt} & 5 \\
\texttt{combi} & 4 \\
\texttt{va} & 4 \\
\texttt{a} & 3 \\
\texttt{go} & 3 \\
\texttt{les} & 3 \\
\texttt{ne} & 3 \\
\texttt{piol} & 3 \\
\texttt{à} & 3 \\
\texttt{ashia} & 2 \\
\texttt{avec} & 2 \\
\texttt{bring} & 2 \\
\bottomrule
\end{tabular}

\end{minipage}
\end{center}

\subsection{Terminal classes}

The grammatical shape of the data, which is what predicts whether a grammar will parse it.

\begin{center}\begin{tabular}{@{}lr@{}}
\toprule
\textbf{Terminal} & \textbf{Count} \\
\midrule
\texttt{NOUN} & 60 \\
\texttt{VERB} & 58 \\
\texttt{DET} & 40 \\
\texttt{PREP} & 25 \\
\texttt{SEP} & 20 \\
\texttt{PRON} & 19 \\
\texttt{ADV} & 14 \\
\texttt{NEG} & 9 \\
\texttt{ADJ} & 6 \\
\texttt{CONJ} & 4 \\
\texttt{INTJ} & 2 \\
\texttt{FORMULA} & 1 \\
\bottomrule
\end{tabular}
\end{center}

\begin{keypoint}[What the variation tells us]
Variation in this data is not noise to be normalised away. The gap between the raw and
canonical tables is a measurement of how unstable written Francanglais orthography is,
and it is precisely the reason the lexer needs a folded fallback tier rather than a single
exact-match dictionary.
\end{keypoint}
